\section{A bounded extension to a single-site update}
\label{sec:single-site}

The balance argument also handles the other single-site endpoint of
the masked-neighbor family.  This extension is included to specify what
the pre-release strengthening check establishes and what it leaves open.

\begin{proposition}[One binary neighbor update]\label{prop:single-site}
\formal{CompositionSingleSiteEndpoint.lean}{single_at_most_two_classes_impossible}
Let $d\ge3$ and let $L=I+E_{0,1}$ act on every binary row/lane word,
where $E_{i,j}x$ has its only possibly nonzero coordinate, $x_j$, at $i$.
For every positive width, there is no complete factorization
$\tau=KBLH$ in which $B,H$ are column-equivariant total maps and $K$
has at most two receiver classes.
\end{proposition}
\begin{outline}
Absorb $B$ into $K$.  We prove adjacent-pair transitivity for
$Q=LTL^{-1}$ and then apply \cref{thm:balanced-pair}.
Put $U_i=I+E_{i,i+1}$, with cyclic indices.  Since $L=U_0=L^{-1}$,
conjugating $\langle Q,TQT^{-1}\rangle$ by $L^{-1}$ gives
$\langle T,A\rangle$, where
\[
 A=U_1U_0U_{-1}U_0.
\]
For $d\ge4$, direct multiplication gives $A=I+N$, with
\[
 N=E_{1,2}+E_{-1,0}+E_{-1,1},\qquad
 N^2=E_{-1,2},\qquad N^3=0.
\]
In characteristic two, $A^2=I+E_{-1,2}$.  Conjugating by rotations
therefore gives every transvection $I+E_{i,i+3}$.

If $d=4$, the directed step-three edges form one cycle.  The elementary
identity
\[
 [I+E_{i,j},I+E_{j,k}]=I+E_{i,k}
 \qquad(i,j,k\text{ distinct})
\]
along simple directed paths yields every $I+E_{i,j}$ with $i\ne j$.
For $d\ge5$, use the already obtained $E=I+E_{0,3}$.
Here $Ae_0=e_0+e_{-1}$, while $e_3^*A^{-1}=e_3^*$, because row three
is unchanged by $A$ and its inverse.  Thus
\[
 (AEA^{-1})E=I+E_{-1,3}.
\]
Rotations also give every step-four transvection.  Within each orbit of
the step-three rotation, a simple directed path from $i$ to $i-3$
gives the reverse step-three transvection by repeated commutators.
If that orbit has two vertices, the reverse edge is already a generator.
Since $d\ge5$, the indices $i,i+4,i+1$ are distinct, and
\[
 [I+E_{i,i+4},I+E_{i+4,i+1}]=I+E_{i,i+1}.
\]
The resulting step-one cycle gives every elementary transvection by
the same chain argument.  Such transvections
act transitively on nonzero binary words: preserve one nonzero coordinate
and add it to clear every other nonzero coordinate, then rotate.

When $d=3$, the coincident indices require a separate calculation.
The same product is
$A(a,b,c)=(a,a+c,a+b+c)$.  It sends $100$ to $111$ and then to $101$.
Their rotations include the three unit vectors, the three weight-two
vectors and the all-one vector, proving transitivity on all seven
nonzero words.  This completes the adjacent-pair argument for every
$d\ge3$.

Hence the partition of $K$ must properly color the cycle.  Odd $d$
is impossible with two colors.  For even $d\ge4$ a proper two-coloring
has period two.  But $L(1,\ldots,1)=(0,1,\ldots,1)$, so no nonzero
column rotation preserves $L\operatorname{Fix}(T)$.  The period
constraint in \cref{thm:period-transport} rules out this last case.
\end{outline}

For clarity, this does not settle the family
$L_m=I+\operatorname{diag}(m)T$ with an arbitrary binary mask $m$.
If $m$ has at least one zero, the weighted shift is nilpotent and
$L_m^{-1}=\sum_{j=0}^{d-1}(\operatorname{diag}(m)T)^j$.
The all-one mask is singular.  A period $r$ of $m$ makes $L_m$ commute
with $T^r$ and gives the class bound of \cref{thm:periodic-first}.
The masks with exactly one zero and exactly one one are covered by
\cref{cor:anchored-all,prop:single-site}, respectively.  The remaining
aperiodic masks, and half-period masks with two final classes, are not
excluded by these arguments.  All statements concern complete algebraic
maps; no native implementation or physical cost bound is asserted.
